Read in Gauge Theory Gravity: the gravitational field $(\hb, \Omega)$ lives on flat spacetime, the ambient connection is a $\Cl(2,4)$ gauge field, and component expressions are labelled readings in an orthonormal frame of the derived metric. Vectors of $\Cl(2,4)$ are $V = \sigma\,\nInf + \mu + \rho\,\nO$ with $Q(V) = 2\sigma\rho + Q(\mu)$, as in Theorem \claimref{synthesis.gtg-tractor-gauge-field-reproduces-tractor-connection}, whose connection is the tractor connection of Fact \claimref{interior.imported-tractor-connection} in the splitting of $g$. Bivectors act on vectors by the commutator product, $(a \wedge b) \times c = (a \wedge b) \cdot c = a\,\ip{b}{c} - b\,\ip{a}{c}$.

\emph{(i)} Let $u = \hb^{-1}(a)$, a spacetime vector, so $\ip{u}{\nInf} = \ip{u}{\nO} = 0$. Since $\Omega(a)$ is a bivector of spacetime vectors and $I_{\rho}$ lies in the plane of $\nInf$, $\nO$, $\Omega(a) \times I_{\rho} = 0$. Next $(\nInf \wedge u) \cdot \nInf = 0$ and $(\nInf \wedge u) \cdot \nO = -u$, so $-(\nInf \wedge u) \times I_{\rho} = \rho u$; $(\nO \wedge S(u)) \cdot \nInf = -S(u)$ and $(\nO \wedge S(u)) \cdot \nO = 0$, so $(\nO \wedge S(u)) \times I_{\rho} = -S(u)$; and $(\nO \wedge \nInf) \cdot \nInf = -\nInf$, $(\nO \wedge \nInf) \cdot \nO = \nO$, so $\beta(a)(\nO \wedge \nInf) \times I_{\rho} = -\beta(a)(\nInf - \rho\,\nO)$. Adding, $A(a) \times I_{\rho} = \rho u - S(u) - \beta(a)(\nInf - \rho\,\nO)$, and $a \cdot \nabla I_{\rho} = (a \cdot \nabla\rho)\,\nO$. The spacetime part of $D_{a}I_{\rho}$ is $\rho u - S(u)$, the $\nInf$ part is $-\beta(a)$ and the $\nO$ part is $a \cdot \nabla\rho + \rho\beta(a)$; they vanish for all $a$ exactly when $S = \rho\,\mathrm{id}$ (the frame is invertible on the invertible region, so $A$ ranges over all vectors), $\beta = 0$ and $\rho$ is constant. No property of $(\hb, \Omega)$ was used: torsion, curvature and matter do not enter. With $S = \rho\,\mathrm{id}$, $\beta = 0$, the connection is $\Omega(a) - (\nInf - \rho\,\nO) \wedge u$, which is the connection $\Omega(a) + \hb^{-1}(a) \wedge V/Q(V)$ of Theorem \claimref{synthesis.curvature-only-tractor-actions-are-macdowell-mansouri}(iv) written in the stored chart: $(\nInf - \rho\,\nO) = W$ of Theorem \claimref{synthesis.gtg-tractor-gauge-field-reproduces-tractor-connection}(v) for $\rho = \Lambda/6$. The norm $Q(I_{\rho}) = 2\rho$ and $\ip{X}{I_{\rho}} = \ip{\nO}{\nInf} = 1$; a connection by bivectors preserves $Q$, so the norm is constant along any curve. In particular, on the field space of connections preserving $I$ the constant $\rho$, hence $\Lambda = 6\rho$, is a fixed parameter. \emph{Gauge.} All of this is in the gauge where $I = \nInf + \rho\,\nO$ in the scale $g$: under the special-conformal rotor $U = \exp(\tfrac12\nO \wedge y)$, which fixes $X = \nO$, the section $I$ acquires a spacetime component and the transformed slot is no longer $\rho\,\mathrm{id}$, so the gauge-invariant content is that $I$ is parallel for $A$ with $\ip{X}{I} = 1$, a statement about the pair $(A, I)$. The check verifies the formula for $A(a) \times I_{\rho}$ exactly at $\rho \in \{\tfrac12, -\tfrac32, 0, \tfrac76\}$ for three frame vectors, three slot values (including $\rho u$) and two dilation values, with a general Lorentz bivector; that the linear map $(S(u), \beta) \mapsto A(a) \times I_{\rho}$ has rank $5$, so the solution is unique; and that a slot $\rho u$ plus any nonzero vector breaks parallelism.

\emph{(ii)} For $D^{T}$, $S = \Sch$ and $\beta = 0$, so by (i) $D^{T}_{a}I_{\rho} = (\rho - \Sch)(u)$ for constant $\rho$, which is Theorem \claimref{synthesis.gtg-tractor-gauge-field-reproduces-tractor-connection}(ii) and, in tractor components, the parallel-section criterion of Theorem \claimref{interior.scale-tractor-leaves-cone}. With the field equation of Theorem \claimref{synthesis.macdowell-mansouri-principle-fixes-the-trace-free-part}(iv), $\Ric + 8\pi G\,T = (\Lambda + 4\pi G\,\tau)g$: tracing, $R = 4\Lambda + 8\pi G\,\tau$, and $\Sch = \tfrac12(\Ric - \tfrac{R}{6}g) = (\tfrac{\Lambda}{6} + \tfrac{4\pi G\tau}{3})g - 4\pi G\,T$, the stored on-shell value. Hence $(\rho - \Sch)(u) = 4\pi G(T - \tfrac{\tau}{3}\,\mathrm{id})(u)$ with $\rho = \Lambda/6$. If $T - \tfrac{\tau}{3}\,\mathrm{id} = 0$ then tracing gives $\tau - \tfrac43\tau = 0$, so $\tau = 0$ and $T = 0$: the map is injective. The difference of the two connections is $\nO \wedge (\Sch - \rho)(u) = -4\pi G\,\nO \wedge (T - \tfrac{\tau}{3}\,\mathrm{id})(u)$. The check computes $\Sch$ from the field equation as matrices in an orthonormal frame for three stress tensors (a general symmetric one, pure dust, pure trace) at $\Lambda \in \{3, -\tfrac32\}$, with the constant $4\pi G$ a rational symbol (the identity is linear in it), confirms the on-shell value and $D^{T}_{a}I_{\rho} - D^{\rho}_{a}I_{\rho} = 4\pi G(T - \tfrac{\tau}{3}\,\mathrm{id})(u)$ with $D^{\rho}_{a}I_{\rho} = 0$ for three frame vectors, and that $T \mapsto T - \tfrac{\tau}{3}\,\mathrm{id}$ has rank $10$ on the ten symmetric tensors.

\emph{(iii)} By (i) the tractor $(\sigma, \mu, \rho')$ is $D^{T}$-parallel exactly when, in the slots of Fact \claimref{interior.imported-tractor-connection}, $\mu = \nabla\sigma$, $\nabla_{a}\nabla_{b}\sigma + g_{ab}\rho' - \Sch_{ab}\sigma = 0$ and $\nabla_{a}\rho' + \Sch_{a}{}^{b}\mu_{b} = 0$; the trace of the middle slot gives $\rho' = \tfrac14(\sigma J - \Delta\sigma)$ with $\Delta = g^{ab}\nabla_{a}\nabla_{b}$, and the norm is $2\sigma\rho' + g^{ab}\nabla_{a}\sigma\nabla_{b}\sigma$. The check builds the metric $g = a(\tau)^{2}\eta$, $\eta = \mathrm{diag}(1, -1, -1, -1)$, its Christoffel symbols and Ricci tensor as exact rational functions of $\tau$, with the sign of the store fixed by its own convention that de Sitter space has $\Ric = \Lambda g$ in signature $(1,3)$ (Theorem \claimref{dark-energy.vacuum-region-registration-condition}(i)): for $a = 1/(H\tau)$ it finds $\Ric = 3H^{2}g$ and the own-scale tractor $(1, 0, \Lambda/6)$ parallel with norm $H^{2} = \Lambda/3$ at $H = 1, \tfrac23$, as Theorem \claimref{interior.scale-tractor-leaves-cone} requires. For $a = \tau^{2}$ (the matter-dominated dust field with $\Lambda = 0$, standard and cited from memory; only the computed curvature is used) it finds $R = 12/\tau^{6}$ and a nonzero trace-free Ricci part, so the field is not Einstein; the own-scale tractor is not parallel and has norm $J/2 = \tau^{-6}$; $\sigma = h\tau^{3}$ gives all three slot equations satisfied with norm $h^{2}$ at $h = 1, 3, -\tfrac25$; $\sigma = \tau^{2} + h\tau^{3}$ gives parallel tractors of constant norm at $h = 0, 1, -\tfrac52$, null at $h = 0$; and $\sigma = \tau, \tau^{4}, \tau^{5}$ violate the middle slot. These are the expected scales: $\sigma^{-2}g = H^{-2}\tau^{-2}\eta$ is de Sitter for $\sigma = H\tau^{3}$ and $\sigma^{-2}g = \eta$ for $\sigma = \tau^{2}$, and the conformally flat dust field has a flat normal tractor connection (its Weyl curvature vanishes). For each of them $\ip{X}{I} = \sigma$ is not constant. Parallel tractors form a vector space and $Q(kI) = k^{2}Q(I)$, so the norm of a transport-parallel tractor is not tied to the physical field unless its scale is the physical one, which by Theorem \claimref{interior.scale-tractor-leaves-cone} happens only on Einstein fields.

\emph{(iv)} In the scale $g$ the splitting of Fact \claimref{interior.imported-tractor-connection} trivialises the tractor bundle as $\R\nInf \oplus TM \oplus \R\nO$, with $\nInf$ and $\nO$ global sections, so $I_{\rho}$ is a global section with $Q(I_{\rho}) = 2\rho$ for every $\rho$; a scale exists on any region where the frame is invertible (the derived metric is one). The orbit map $\xi \mapsto \xi \times I_{\rho}$ from the fifteen bivectors to vectors has image in $I_{\rho}^{\perp}$ (since $\ip{\xi \times I}{I} = 0$ for a bivector $\xi$) and the check finds rank $5 = \dim I_{\rho}^{\perp}$ for $\rho \in \{\tfrac12, 5, -\tfrac13, 0\}$, so the stabiliser has dimension $10$ (Fact \claimref{dark-energy.imported-frozen-point-geometry}), and in the basis $e_{0}, \ldots, e_{3}, \nInf - \rho\,\nO$ of $I_{\rho}^{\perp}$ the Gram matrix has signature $(1,4)$ at $\rho = \tfrac12, 5$ and $(2,3)$ at $\rho = -\tfrac13$ (the last vector has norm $-2\rho$). Hence the gauge group acts locally transitively on each level set of the norm, and a section of fixed norm is gauge equivalent, on a contractible region, to the constant one; what a dynamical frozen point with the connection unrestricted gives is proved, within a family of functionals, in Theorem \claimref{synthesis.frozen-point-reduces-the-connection-only-when-infinitely-stiff}. \emph{(Topology, argued.)} Writing a vector of $\R^{2,4}$ as $(x_{+}, x_{-}) \in \R^{2} \times \R^{4}$ with $Q = |x_{+}|^{2} - |x_{-}|^{2}$, the level set $Q = c > 0$ is $\{x_{+} = \sqrt{c + |x_{-}|^{2}}\,\theta\}$, a circle times $\R^{4}$, and $Q = c < 0$ is a three-sphere times $\R^{2}$; the obstructions to a section of a bundle with such fibres lie in degrees $2$ and $4$, but they vanish here because the section $I_{\rho}$ is exhibited. That the sign of $\Lambda$ is otherwise free is Theorem \claimref{odd-extensions.transport-constrains-frozen-point-only-through}(iii)-(iv); the check confirms that $Q(I_{\rho})$ takes each sign.

\emph{Reading.} Clause (iii) of Theorem \claimref{dark-energy.vacuum-region-registration-condition}, in the paper's wording, names its condition without specifying the connection (its current wording names the connection preserving $I$ and the physical scale, Premise \claimref{synthesis.field-space-premise-parallel-frozen-point}); read literally, (i) meets it for every field through $A_{\rho}$. With the series' transport it fails on matter regions in the physical scale by (ii), which is on shell for the frozen-point functional with the imported matter coupling, so it uses Einstein's equations with matter as the field equations of that functional. By (i) the only connection of the soldered tractor shape that keeps $I_{\rho}$ parallel is $A_{\rho}$; hence adopting the field space of Theorem \claimref{synthesis.curvature-only-tractor-actions-are-macdowell-mansouri}(iv) as the gravitational connection implies the condition, and conversely the condition read kinematically (the connection that carries the frozen point is, for every configuration, a soldered connection keeping it parallel in the physical scale) is that field space. The converse fails for the reading on solutions only, since on the unreduced space a functional can make the slot stationary away from $\rho\,\mathrm{id}$, or stationary at $\rho\,\mathrm{id}$ with only the Euler density left (Theorem \claimref{synthesis.frozen-point-reduces-the-connection-only-when-infinitely-stiff}). Clause (ii) of Theorem \claimref{dark-energy.vacuum-region-registration-condition} shows why the physical scale is the relevant reading: the norm must be the $\Lambda$ entering the physical field equations, and (iii) here shows that without it the condition is met by transport on conformally flat matter regions with any norm. On the reduced space $\rho$ is constant by (i), so the field equation is Einstein's with a constant $\Lambda$ (Theorem \claimref{synthesis.macdowell-mansouri-principle-fixes-the-trace-free-part}(iv)) and $w = -1$ there is the $\Lambda$CDM equation of state. The class of connections is the soldered tractor shape (the frame slot fixed as the soldering form, the rotation block that of the field); connections outside it are not treated.
